% Generated draft: requires mathematical and source review.
\documentclass[11pt]{article}
\usepackage[margin=25mm]{geometry}
\usepackage{fontspec}
\setmainfont{Latin Modern Roman}
\usepackage{amsmath,amssymb,mathtools,unicode-math}
\setmathfont{Latin Modern Math}
\usepackage{enumitem,needspace,xurl,hyperref}
\setlength{\parindent}{0pt}
\setlength{\parskip}{5pt}
\setlength{\emergencystretch}{3em}
\newcommand{\N}{\mathbb{N}}
\newcommand{\Z}{\mathbb{Z}}
\newcommand{\Q}{\mathbb{Q}}
\newcommand{\R}{\mathbb{R}}
\newcommand{\C}{\mathbb{C}}
\newcommand{\F}{\mathbb{F}}
\newcommand{\A}{\mathbb{A}}
\newcommand{\PP}{\mathbb{P}}
\newcommand{\E}{\mathbb{E}}
\newcommand{\eps}{\varepsilon}
\newcommand{\dd}{\,\mathrm{d}}
\begin{document}
\section*{UnsolvedMath ID: 30005676}
\textbf{Source problem code:} catalogue entry
\subsection{Definitions and mathematical statement}
Let \(\mathcal M_{\mathrm{ell}}\) denote the moduli stack of smooth elliptic curves, and let \(\mathcal O^{\mathrm{top}}\) denote the topological modular-form construction referred to in the supplied report excerpt. The supplied assessment describes an established isogeny-indexed construction over \(\mathcal M_{\mathrm{ell}}\); the excerpt itself does not include its theorem, hypotheses, or definition of the indexing category \(\mathrm{Isog}\). Write \(\mathcal G\) for the intended moduli stack of generalized elliptic curves. The supplied assessment describes this as a compactified target, but does not specify the precise stack, base, or permitted degenerations.

Let \(\mathbb S\) be the sphere spectrum. Write \(L_{\mathrm{TMF}}\mathbb S\) for its Bousfield localization with respect to the spectrum \(\mathrm{TMF}\). The notation \(\mathrm{Tmf}\) in the source is retained separately from \(\mathrm{TMF}\), as required by its question about extending the construction to generalized elliptic curves.

\begin{enumerate}
\item Can the report's theorem and its topological modular-form operations be extended from \(\mathcal M_{\mathrm{ell}}\) to \(\mathcal G\), compatibly with restriction to smooth elliptic curves and with the isogeny operations in that theorem? Can this give an approach to \(\mathrm{Tmf}\) analogous to the approach to \(\mathrm{TMF}\) attributed in the excerpt to Lurie and its reference~[5]? The omitted theorem and unidentified reference~[5] are necessary to specify the exact extension problem.
\item After specifying the report's category \(\mathrm{Isog}\) and the associated diagram of spectra, identify its limit
\[
X:=\operatorname*{lim}_{\mathrm{Isog}}\mathcal O^{\mathrm{top}}.
\]
Here the limit is understood in the homotopical sense, as a limit in the appropriate infinity-category of spectra or structured ring spectra. The supplied record does not specify the diagram's variance, its objects and morphisms, its transition maps, or the passage from sheaves to spectra, so this display names the intended target without determining a unique spectrum.
\item For the proposed comparison from the \(\mathrm{TMF}\)-local sphere,
\[
u\colon L_{\mathrm{TMF}}\mathbb S\longrightarrow X,
\]
determine in what precise sense it approximates that sphere, and whether the limit diagram produces an integral or near-integral resolution of \(L_{\mathrm{TMF}}\mathbb S\). The direction of this proposed map follows the supplied reconstructed statement; its construction is not given. In particular, neither a factorization of a unit map through \(L_{\mathrm{TMF}}\mathbb S\) nor the necessary locality of the target is established by the supplied excerpt. No meaning of ``good approximation,'' ``near-integral,'' or ``resolution'' is specified, including any allowed primes to invert, diagrammatic form, or convergence requirement.
\end{enumerate}

These are qualified research questions, not a claimed extension theorem or equivalence. The missing construction data and success criteria prevent a fully determined mathematical conjecture from being extracted from this record alone.
\subsection{Short English statement}
Extend the report's topological modular-form construction to generalized elliptic curves, identify its isogeny-indexed homotopy limit, and investigate whether the proposed comparison from the \(\mathrm{TMF}\)-local sphere provides an integral or near-integral resolution. The excerpt omits the preceding theorem and essential definitions, so the exact extension, comparison map, and resolution criteria remain unresolved. The supplied open status is dated curation, not an independently verified conclusion.
\subsection{Sources}
\begin{itemize}
\item Local UnsolvedMath registry: supplied complete record, permanent ID \(30005676\), source problem code \texttt{catalogue entry}, title ``Topological Modular Forms over Generalized Elliptic Curves.'' No individual proposer is supplied; the recorded proposed year is \(2023\).
\item Supplied report citation: \emph{Homotopy Theory (2024)}, Oberwolfach Reports, pp.~1893--1952, DOI \texttt{10.4171/owr/2023/34}. \url{https://doi.org/10.4171/owr/2023/34}. This is the supplied authoritative report reference; its contents have not been independently consulted for this draft.
\item Supplied literature source: ``Davies on Hecke operators.'' \url{https://arxiv.org/abs/2212.06208}.
\item Supplied literature source: ``Davies on dualisable Tmf.'' \url{https://arxiv.org/abs/2104.13407}.
\item Literature assessment dated exactly \texttt{2026-08-22}: the supplied curation labels the problem open, describes the isogeny-indexed construction over the elliptic-curve moduli stack as established, and describes the compactified extension and corresponding integral limit description of the \(\mathrm{TMF}\)-local sphere as remaining questions. This is supplied dated curation, not independently verified literature assessment. The background explicitly requests independent verification of time-sensitive status and cited claims.
\item The registry labels the reconstructed statement \texttt{corrected\_verified} and says it was checked against the official report. These are supplied verification claims, not verification performed for this draft. The excerpt attributes an approach to \(\mathrm{TMF}\) to Lurie through reference~[5], whose bibliographic identity is not supplied.
\end{itemize}
\end{document}
